% Fictional teaching example. Replace all example identities and URLs before
% publishing a real edition. The project URL below is intentionally a placeholder.
\ResetPaperMetadata
\PaperTitle{Cross-Platform AI Trend Radar}
\PaperInstitution{Faculty of Media, Hochschule D\"usseldorf}
\PaperEmail{}
\PaperAbstract{Trend lists from a single platform can mistake temporary attention for a durable topic and can miss a topic expressed with different words elsewhere. We design a cross-platform AI trend radar that normalizes weekly signals from discussion, search, and news streams; groups related phrases with sentence embeddings; and ranks clusters by growth, activity, and source diversity. A dashboard exposes the score components and offers a source-grounded summary draft. In a small, explicitly illustrative ten-topic exercise, the combined rank places four relevant topics in its top five, compared with three for a volume-only baseline. The example demonstrates an inspectable evaluation workflow, while its synthetic data cannot establish forecasting performance.}
\PaperKeywords{trend detection, semantic clustering, cross-platform signals, evaluation}
\ProjectRepository{https://github.com/your-account/cross-platform-ai-trend-radar}

% Each person appears in this order on the paper and in About the Authors.
% Photo fields are empty so the neutral fallback is exercised in this example.
\AddAuthor
  {Alex Rivera}
  {B.Sc. Data Science}
  {Alex Rivera is a fictional Data Science student at Hochschule D\"usseldorf. In the illustrative Cross-Platform AI Trend Radar project, Alex focused on source normalization and on making ranking decisions traceable to weekly counts. Their interests include information retrieval, evaluation design, and responsible use of language models. This profile is sample content for the proceedings template and does not describe a real person.}
  {}
  {}
  {}
  {https://example.com} % Reserved example domain, not a personal portfolio.
\AddAuthor
  {Samira Becker}
  {B.Sc. Data Science}
  {Samira Becker is a fictional Data Science student at Hochschule D\"usseldorf. In the illustrative Cross-Platform AI Trend Radar project, Samira focused on semantic grouping and on inspecting cases where different topics were merged. Her interests include natural language processing, transparent evaluation, and human review of generated summaries. This profile is sample content for the proceedings template and does not describe a real person.}
  {}
  {}
  {}
  {}

